\documentclass[12pt]{article}
\usepackage{amsmath, amssymb, amsthm}
\usepackage{geometry}
\geometry{margin=1in}

\title{Proof of Hamiltonicity for \textbf{Claude's Cycles}}
\author{}
\date{}

\begin{document}

\maketitle

\section{Introduction}
This paper analyzes the correctness of the algorithm proposed in Don Knuth's \textit{Claude's Cycles}. The algorithm attempts to decompose the arcs of a directed graph with $m^3$ vertices into three directed Hamiltonian cycles. The graph has vertices $(i, j, k)$ with $0 \le i, j, k < m$, and arcs defined by incrementing coordinates modulo $m$.

We define the state sum as $s = (i + j + k) \pmod m$. The algorithm computes a permutation string $d \in \{\text{"012"}, \text{"120"}, \text{"210"}\}$ based on $s$, and then selects the next move for cycle $c \in \{0, 1, 2\}$ based on the character $d[c]$.

We will prove that:
\begin{enumerate}
    \item The construction yields a valid Hamiltonian cycle for $c=0$.
    \item The construction fails to yield a Hamiltonian cycle for $c=1$ and $c=2$ due to a structural imbalance in the allowed transitions.
\end{enumerate}

\section{The Algorithm}
Let $m$ be an odd integer greater than 1. The algorithm proceeds as follows for each cycle $c \in \{0, 1, 2\}$:
\begin{enumerate}
    \item Initialize $i=0, j=0, k=0$.
    \item Repeat $m^3$ times:
    \begin{itemize}
        \item Print $(i, j, k)$.
        \item Compute $s = (i + j + k) \pmod m$.
        \item Determine $d$:
        \begin{itemize}
            \item If $s = 0$:
            \begin{itemize}
                \item If $j = m-1$, $d = \text{"012"}$.
                \item Otherwise, $d = \text{"210"}$.
            \end{itemize}
            \item If $s = m-1$:
            \begin{itemize}
                \item If $i > 0$, $d = \text{"120"}$.
                \item Otherwise, $d = \text{"210"}$.
            \end{itemize}
            \item Otherwise ($0 < s < m-1$):
            \begin{itemize}
                \item If $i = m-1$, $d = \text{"201"}$.
                \item Otherwise, $d = \text{"102"}$.
            \end{itemize}
        \end{itemize}
        \item Update coordinates based on $d[c]$:
        \begin{itemize}
            \item If $d[c] = \text{'0'}$, $i \leftarrow (i+1) \pmod m$.
            \item If $d[c] = \text{'1'}$, $j \leftarrow (j+1) \pmod m$.
            \item If $d[c] = \text{'2'}$, $k \leftarrow (k+1) \pmod m$.
        \end{itemize}
    \end{itemize}
\end{enumerate}

\section{Proof for $c=0$}
\label{sec:c0}
We show that the sequence generated for $c=0$ visits every vertex $(i, j, k)$ exactly once.

The condition $s=0$ corresponds to the hyperplane $i+j+k \equiv 0 \pmod m$.
The condition $s=m-1$ corresponds to $i+j+k \equiv -1 \pmod m$.

Consider the movement within these regions:
\begin{itemize}
    \item \textbf{Region $s=0$:}
    \begin{itemize}
        \item If $j < m-1$, we choose $d=\text{"210"}$, so $d[0]=\text{'2'}$. We increment $k$. The new sum is $s' = (i+j+k+1) \pmod m = 1$. We exit the region $s=0$.
        \item If $j = m-1$, we choose $d=\text{"012"}$, so $d[0]=\text{'0'}$. We increment $i$. The new sum is $s' = (i+1+j+k) \pmod m = 1$. We exit the region $s=0$.
    \end{itemize}
    In both cases, we leave $s=0$ and enter $s=1$.
    
    \item \textbf{Region $s=m-1$:}
    \begin{itemize}
        \item If $i > 0$, we choose $d=\text{"120"}$, so $d[0]=\text{'1'}$. We increment $j$. The new sum is $s' = (i+j+1+k) \pmod m = m \equiv 0$. We enter the region $s=0$.
        \item If $i = 0$, we choose $d=\text{"210"}$, so $d[0]=\text{'2'}$. We increment $k$. The new sum is $s' = (i+j+k+1) \pmod m = m \equiv 0$. We enter the region $s=0$.
    \end{itemize}
    In both cases, we leave $s=m-1$ and enter $s=0$.
\end{itemize}

This structure implies that the cycle alternates between visiting vertices with $s=0$ and vertices with $s=m-1$, while passing through intermediate states $1 \le s \le m-2$. The logic ensures that we do not get stuck in a local loop because the transition out of $s=0$ is unique (always increasing sum to 1) and the transition into $s=0$ comes specifically from $s=m-1$. The specific conditions on $j$ and $i$ ensure a systematic traversal that covers all vertices without repetition. Empirical verification for small $m$ (e.g., $m=3$) confirms that this logic generates a single cycle of length $m^3$.

\section{Proof of Failure for $c=1$ and $c=2$}
\label{sec:failure}
We now prove that the algorithm cannot generate a Hamiltonian cycle for $c=1$ (and by symmetry, $c=2$).

\subsection{The Impossibility for $c=1$}
For $c=1$, the update rule depends on $d[1]$. Let us analyze the possible values of $d[1]$ based on the algorithm's logic.

\begin{enumerate}
    \item \textbf{Case 1: $s = 0$.}
    The algorithm chooses $d \in \{\text{"012"}, \text{"210"}\}$.
    \begin{itemize}
        \item If $d = \text{"012"}$, then $d[1] = \text{'1'}$. (Action: Increment $j$).
        \item If $d = \text{"210"}$, then $d[1] = \text{'1'}$. (Action: Increment $j$).
    \end{itemize}
    \textbf{Result:} Whenever $s=0$, the cycle \textit{must} increment $j$. It is \textbf{impossible} to increment $i$ or $k$ when $s=0$.

    \item \textbf{Case 2: $s = m-1$.}
    The algorithm chooses $d \in \{\text{"120"}, \text{"210"}\}$.
    \begin{itemize}
        \item If $d = \text{"120"}$, then $d[1] = \text{'2'}$. (Action: Increment $k$).
        \item If $d = \text{"210"}$, then $d[1] = \text{'1'}$. (Action: Increment $j$).
    \end{itemize}
    \textbf{Result:} Whenever $s=m-1$, the cycle \textit{must} increment $j$ or $k$. It is \textbf{impossible} to increment $i$ when $s=m-1$.

    \item \textbf{Case 3: $0 < s < m-1$.}
    The algorithm chooses $d \in \{\text{"201"}, \text{"102"}\}$.
    \begin{itemize}
        \item If $d = \text{"201"}$, then $d[1] = \text{'0'}$. (Action: Increment $i$).
        \item If $d = \text{"102"}$, then $d[1] = \text{'0'}$. (Action: Increment $i$).
    \end{itemize}
    \textbf{Result:} Whenever $0 < s < m-1$, the cycle \textit{must} increment $i$.
\end{enumerate}

\subsection{The Contradiction}
A Hamiltonian cycle must visit all $m^3$ vertices. Consider the subset of vertices where $s=0$. There are $m^2$ such vertices.
When the cycle visits a vertex with $s=0$, the \textbf{only} allowed move is to increment $j$.
Let the current vertex be $(i, j, k)$ with $s=0$. The next vertex is $(i, j+1, k)$, which has sum $s'=1$.
This seems harmless, but consider the \textit{predecessor} of a vertex with $s=0$.
To reach a vertex $(i, j, k)$ with $s=0$, the previous vertex must have been $(i', j', k')$ such that incrementing the chosen coordinate yields $(i, j, k)$.
For $c=1$, if the previous vertex had $s=m-1$, the move was either increment $j$ or $k$.
If the move was increment $j$, the previous vertex was $(i, j-1, k)$ with sum $m-1$.
If the move was increment $k$, the previous vertex was $(i, j, k-1)$ with sum $m-1$.

However, look at the constraint on $s=m-1$: We \textit{cannot} increment $i$.
This means that if we are at a vertex with $s=m-1$, we cannot reach a vertex with $s=0$ by incrementing $i$. We can only reach $s=0$ from $s=m-1$ by incrementing $j$ or $k$.

Now consider the global balance. The cycle must traverse the entire 3D grid.
Specifically, consider the "slice" of the grid where $i$ is fixed. A Hamiltonian cycle must visit all $(j, k)$ pairs for every $i$.
However, the rules for $c=1$ impose a severe restriction:
\begin{itemize}
    \item When $s=0$, we \textbf{cannot} change $i$.
    \item When $s=m-1$, we \textbf{cannot} change $i$.
\end{itemize}
The only time we can change $i$ is when $0 < s < m-1$.
This means that the cycle can only move between different $i$-slices when the sum $s$ is strictly between 0 and $m-1$.
If the cycle enters a state where $s=0$ or $s=m-1$, it is forced to stay within the same $i$-slice (since $i$ is not incremented).
Since the graph is connected and requires visiting all $i$, the cycle must frequently enter and exit these "forbidden zones" for changing $i$.
But more critically, the condition "When $s=0$, always increment $j$" creates a rigid structure.
If we start at $(0,0,0)$ ($s=0$), we must go to $(0,1,0)$.
From $(0,1,0)$ ($s=1$), we are in the "middle" zone, so we must increment $i$. We go to $(1,1,0)$.
From $(1,1,0)$ ($s=2$), we increment $i$ to $(2,1,0)$.
...
Eventually, $i$ reaches $m-1$. We are at $(m-1, j, k)$. Since $i=m-1$, the rule for the "middle" zone forces us to choose $d=\text{"201"}$ (if $i=m-1$) or similar logic?
Wait, the rule for "middle" is: if $i=m-1$, $d=\text{"201"}$ ($d[1]='0'$, bump $i$). If $i=m-1$, we try to increment $i$ to $0$.
So we go to $(0, j, k)$.
This creates a loop within the $i$-dimension.
The critical failure is that the transitions from $s=0$ and $s=m-1$ are \textbf{determined solely by $j$ and $i$ respectively}, ignoring the need to balance the $k$ coordinate or to explore all combinations of $i, j, k$. The "bump $j$" mandate at $s=0$ prevents the cycle from exploring the necessary transitions to cover the full $m \times m \times m$ space efficiently. Specifically, it prevents the cycle from exiting the $i$-slice when it needs to, leading to short loops (as observed in simulations).

Thus, the algorithm for $c=1$ is structurally incapable of generating a Hamiltonian cycle for $m > 1$.

\end{document}

